% =============================================================
%  Section 08: 结论与展望
% =============================================================
\section{结论与展望}

\begin{frame}{结论：本文回答了什么}
  \begin{block}{结论 1}
    纯文本 Baseline 在当前实验设定下最稳定；\hl{异构接入并未自动带来收益}。
  \end{block}
  \begin{block}{结论 2}
    Adaptive PRF 的主要价值在于\hl{提升多跳证据覆盖率}，尤其对 bridge 类问题有效。
  \end{block}
  \begin{block}{结论 3}
    当前启发式 CoVe 无法同时兼顾\hl{低误杀}与\hl{低漏放}，尚不足以直接用于高可靠拒答。
  \end{block}
  \begin{alertblock}{核心认识}
    RAG 不是简单的“模块堆叠”问题，而是\hl{效果、成本、安全性三者的系统权衡问题}。
  \end{alertblock}
\end{frame}

\begin{frame}{未来工作}
  \begin{columns}[T,onlytextwidth]
    \column{0.33\textwidth}
    \begin{block}{验证闭环}
      将 CoVe 从“拦截器”升级为“反馈节点”，失败后回写检索与再生成。
    \end{block}

    \column{0.33\textwidth}
    \begin{block}{学习式预算}
      用成本感知策略替代静态阈值，显式建模覆盖收益与 Token/延迟代价。
    \end{block}

    \column{0.33\textwidth}
    \begin{block}{真实异构数据}
      接入真实表格库与知识图谱，替代当前伴生构造，验证结论的外部有效性。
    \end{block}
  \end{columns}

  \vspace{0.6em}
  \begin{exampleblock}{一句话总结}
    本文的价值在于给出一套\hl{可复现、可诊断、可继续扩展}的 RAG 研究基线，而不是给出一个“已经完美”的系统。
  \end{exampleblock}
\end{frame}
